\section{Experiments}
\label{sec:exp}

\subsection{Experimental Setup}

Hestia is trained on our processed Objaverse~\cite{deitke2024objaverse, Deitke_2023_CVPR} split to showcase its full capability and on our Houses3K~\cite{peralta2020next} split denoted as Hestia-H3K for fair comparison. 
We use NVIDIA IsaacLab~\cite{mittal2023orbit} to randomly simulate 256 scenes in parallel, with each object scaled up to 8 meters and placed in a 20$\times$20$\times$20m scene.  
Objects are placed at the origin and the four corners for benchmarking.
%Objects are randomly translated during training and placed at the origin and the four corners for benchmarking.
An RGB-D camera is ahead of the Crazyflie drone, which starts from a random collision-free position oriented toward the object center.
%Unless otherwise specified, we benchmark all methods on three datasets~\cite{deitke2024objaverse, Deitke_2023_CVPR, wu2023omniobject3d, peralta2020next}, resulting in 400 diverse shapes for evaluation, which ensures a comprehensive and fair comparison across methods.
See~\cref{asec:datasets,asec:train_det} for more details.

\subsection{Overall Performance}

This section addresses three questions: \underline{Q1}: Is Hestia’s improvement marginal? \underline{Q2}: Does the method outperform prior works~\cite{jin2023neu, guedon2023macarons, peralta2020next, chen2024gennbv} with and without large-scale training~\cite{deitke2024objaverse, Deitke_2023_CVPR}? \underline{Q3}: Does large-scale training further improve performance? 
To answer these questions, we benchmark on three datasets~\cite{deitke2024objaverse, Deitke_2023_CVPR, wu2023omniobject3d, peralta2020next} comprising 400 diverse shapes, ensuring a comprehensive and fair comparison across methods for the point cloud reconstruction task.
Given the large-scale test set, we select three generalizable baselines~\cite{jin2023neu, chen2024gennbv, peralta2020next} that do not require test-time optimization, along with one online-learning approach~\cite{guedon2023macarons} for benchmarking.
We do not include 3DGS-based online-learning methods~\cite{wilson2025pop, li2025activesplat, khass2025active} as baselines due to differences in data modality, nor methods~\cite{li2025nextbestpath, chen2025gleam} that target non-object-centric scenes with fewer degrees of freedom.


\cref{tab:avg_bench} shows that Hestia not only outperforms prior work on all three datasets, but also achieves at least 4\% and 6\% gains in coverage ratio (CR) and area under the coverage ratio curve (AUC), respectively, while reducing chamfer distance (CD) by 50\% compared to others.
Hence, this answers \underline{Q1}, showing that Hestia’s improvement is not marginal.
Hestia-H3K trained on a smaller, less diverse dataset (Houses3K) still outperforms prior work, demonstrating that the improvement comes not only from large-scale diverse training but also from the proposed designs, thus answering \underline{Q2}.
On both OmniObject3D and Objaverse, Hestia surpasses Hestia-H3K, and even on Hestia-H3K’s own in-distribution set (Houses3K), it achieves slightly better CD, indicating that training on a larger and more diverse dataset provides additional benefits, thus answering \underline{Q3}.

\begin{table*}[thb!]
\centering
\setlength{\tabcolsep}{3.1pt}
\begin{tabular}{cl|
  ccc|
  ccc|
  ccc|
  ccc|c}
\multirow{2}{*}{Train Data} &
\multirow{2}{*}{Method} &
\multicolumn{3}{c|}{OmniObject3D} &
\multicolumn{3}{c|}{Objaverse} &
\multicolumn{3}{c|}{Houses3K} &
\multicolumn{3}{c|}{Overall} &
\multirow{2}{*}{FPS$\uparrow$} \\
& & CR$\uparrow$ & CD$\downarrow$ & AUC$\uparrow$
  & CR$\uparrow$ & CD$\downarrow$ & AUC$\uparrow$
  & CR$\uparrow$ & CD$\downarrow$ & AUC$\uparrow$
  & CR$\uparrow$ & CD$\downarrow$ & AUC$\uparrow$ & \\
    \specialrule{.1em}{.1em}{.1em}
  DTU &  NeU-NBV~\cite{jin2023neu} & 77 & 48 & 73 & 79 & 33 & 75 & 78 & 38 & 75 & 78 & 40 & 74 & $<$1 \\
    \specialrule{.1em}{.1em}{.1em}
    \NA & MACARONS~\cite{guedon2023macarons} & 89 & 17 & 80 & 85 & 24 & 76 & 85 & 26 & 77 & 86 & 22 & 78 & $<$1 \\
    \specialrule{.1em}{.1em}{.1em}
  \multirow{4}{*}{Houses3K} &  ScanRL~\cite{peralta2020next}& 80 & 32 & 75 & 81 & 27 & 78 & 81 & 32 & 76 & 81 & 30 & 76 & 15 \\
    & GenNBV~\cite{chen2024gennbv} & 93 & 12 & \cellcolor{lightyellow}{87} & 91 & 14 & 86 & 92 & 14 & 87 & 92 & 13 & 87 & \cellcolor{lightgreen}{27} \\
    & GenNBV (Rep.)~\cite{chen2024gennbv}& 93 & 11 & \cellcolor{lightyellow}{87} & 92 & 13 & 86 & \cellcolor{lightyellow}{94} & 12 & 89 & 93 & 12 & 87 & \cellcolor{lightgreen}{27} \\
    & Hestia-H3K (Ours) & \cellcolor{lightyellow}{96} & \cellcolor{lightyellow}{6} & \cellcolor{lightgreen}{93}
 & \cellcolor{lightyellow}{95} & \cellcolor{lightyellow}{8} & \cellcolor{lightyellow}{91} & \cellcolor{lightgreen}{97} & \cellcolor{lightyellow}{7} & \cellcolor{lightgreen}{96} & \cellcolor{lightyellow}{96} & \cellcolor{lightyellow}{7} & \cellcolor{lightyellow}{92} & \cellcolor{lightyellow}{25} \\
    \specialrule{.1em}{.1em}{.1em}
    Objaverse & Hestia (Ours)& \cellcolor{lightgreen}{97} & \cellcolor{lightgreen}{4} & \cellcolor{lightgreen}{93} & \cellcolor{lightgreen}{96} & \cellcolor{lightgreen}{7} & \cellcolor{lightgreen}{92} & \cellcolor{lightgreen}{97} & \cellcolor{lightgreen}{6} & \cellcolor{lightyellow}{94} & \cellcolor{lightgreen}{97} & \cellcolor{lightgreen}{6} & \cellcolor{lightgreen}{93} & \cellcolor{lightyellow}{25} \\
    \specialrule{.1em}{.1em}{.1em}
  \end{tabular}
\caption{\textbf{Overall performance on OmniObject3D, Objaverse, and Houses3K with 30 images per object.} Results are reported as mean CR (\%), CD (cm), and AUC (\%) over five object center positions. Hestia and Hestia-H3K outperform prior approaches by at least 4\% and 3\% in CR and by 6\% and 5\% in AUC, respectively, while reducing CD by nearly 50\%. Interestingly, Hestia achieves slightly better CD than Hestia-H3K on in-distribution data (Houses3K).}
\label{tab:avg_bench}
\end{table*}



\subsection{Qualitative Comparisons}
\label{sec:qua}

\begin{figure*}[tbh!]
  \centering
  \includegraphics[width=0.82\linewidth]{figures/overall_vis.png}
  \caption{\textbf{Point cloud reconstruction on three datasets.} Hestia’s reconstructions are visibly better than those of prior approaches.}
  \label{fig:vis_all}
\end{figure*}

% \begin{figure*}[tbh!]
%   \centering
%   \includegraphics[width=\linewidth]{figures/overall_vis.png}
%   \caption{\textbf{Point cloud reconstruction on three datasets.} Hestia’s reconstructions are visibly better than those of prior approaches.}
%   \label{fig:vis_all}
% \end{figure*}

This section highlights that Hestia’s improvement also extends to visualization in the point cloud reconstruction task. 
\cref{fig:vis_all} shows that Hestia produces more comprehensive point clouds than prior work across diverse object shapes. 
Specifically, prior methods fail to reconstruct parts of the teddy bear and anime figurine from OmniObject3D~\cite{wu2023omniobject3d}, the underside of the stair and the cactus's hat and hands from Objaverse~\cite{deitke2024objaverse, Deitke_2023_CVPR}, and self-occluded structures such as the roof soffit or window from Houses3K~\cite{peralta2020next}.
In addition, Hestia performs well on the complex scenes (see~\cref{fig:vis_all2}).
This improvement is largely attributed to our design, which incorporates a hierarchical structure that better identifies missing parts of objects and models voxels as cubes rather than points, thereby preserving geometric details.
More qualitative results (see~\cref{asec:vis_res,asec:limit}), including failure cases, are provided in the supplementary material, and all reconstruction results are included in the supplementary video. 

\subsection{Translation Robustness}

\begin{figure*}[tbh!]
  \centering
  \includegraphics[width=0.82\linewidth]{figures/all_shift.png}
  \caption{\textbf{Point cloud reconstruction on three datasets.} Hestia’s reconstructions are visibly better than those of prior approaches.}
  \label{fig:vis_shift}
\end{figure*}

This section demonstrates that Hestia maintains robustness when objects are placed at different positions within the scenes. 
\cref{atab:base_bench_omni,atab:base_bench_obj,atab:base_bench_h3k} provide detailed results across different object placement settings.
Hestia exhibits less performance fluctuation, outperforming prior methods on all three datasets.
The qualitative results (see~\cref{fig:vis_shift,asec:novel}) also show that Hestia's reconstruction is more robust.
These results highlight the effectiveness of the hierarchical structure, which first predicts the look-at point and then determines the capture destination.

% \begin{table*}[thb!]
%   \centering
%   \setlength{\tabcolsep}{3.1pt}
% \begin{tabular}{l|
%   ccc|ccc|ccc|ccc|ccc}
%      & \multicolumn{15}{c}{OmniObject3D Test} \\
%     \specialrule{.1em}{.1em}{.1em}
%      \multirow{2}{*}{Method} 
%      & \multicolumn{3}{c|}{(0, 0)} 
%      & \multicolumn{3}{c|}{(4, 4)} 
%      & \multicolumn{3}{c|}{(4, -4)} 
%      & \multicolumn{3}{c|}{(-4, 4)} 
%      & \multicolumn{3}{c}{(-4, -4)} \\
%      & CR$\uparrow$ & CD$\downarrow$ & AUC$\uparrow$
%      & CR$\uparrow$ & CD$\downarrow$ & AUC$\uparrow$
%      & CR$\uparrow$ & CD$\downarrow$ & AUC$\uparrow$
%      & CR$\uparrow$ & CD$\downarrow$ & AUC$\uparrow$
%      & CR$\uparrow$ & CD$\downarrow$ & AUC$\uparrow$ \\
%     \specialrule{.1em}{.1em}{.1em}
%      NeU-NBV~\cite{jin2023neu} & 88 & 24 & 83 & 73 & 51 & 70 & 75 & 75 & 70 & 72 & 53 & 69 & 76 & 38 & 71 \\
%     \specialrule{.1em}{.1em}{.1em}
%      ScanRL~\cite{peralta2020next} & 87 & 21 & 80 & 79 & 37 & 76 & 84 & 21 & 78 & 72 & 47 & 71 & 76 & 35 & 72 \\
%     MACARONS~\cite{guedon2023macarons} & 86 & 19 & 75 & 93 & 13 & 86 & 91 & 14 & 83 & \cellcolor{lightyellow}{94} & \cellcolor{lightyellow}{10} & 88 & 80 & 29 & 69 \\
%     GenNBV~\cite{chen2024gennbv} & 92 & 12 & \cellcolor{lightyellow}{88} & 92 & 14 & 86 & 91 & 15 & 83 & \cellcolor{lightyellow}{94} & \cellcolor{lightyellow}{10} & \cellcolor{lightyellow}{89} & \cellcolor{lightyellow}{94} & \cellcolor{lightyellow}{9} & \cellcolor{lightyellow}{88} \\
%     GenNBV (Rep.)~\cite{chen2024gennbv} & \cellcolor{lightyellow}{93} & \cellcolor{lightyellow}{10} & 87 & \cellcolor{lightyellow}{94} & \cellcolor{lightyellow}{9} & \cellcolor{lightyellow}{89} & \cellcolor{lightyellow}{92} & \cellcolor{lightyellow}{12} & \cellcolor{lightyellow}{85} & \cellcolor{lightyellow}{94} & \cellcolor{lightyellow}{10} & 88 & 92 & 11 & 85 \\
%     \specialrule{.1em}{.1em}{.1em}
%     Hestia (Ours) & \cellcolor{lightgreen}{97} & \cellcolor{lightgreen}{4} & \cellcolor{lightgreen}{94} 
%            & \cellcolor{lightgreen}{97} & \cellcolor{lightgreen}{4} & \cellcolor{lightgreen}{93} 
%            & \cellcolor{lightgreen}{97} & \cellcolor{lightgreen}{5} & \cellcolor{lightgreen}{92} 
%            & \cellcolor{lightgreen}{97} & \cellcolor{lightgreen}{4} & \cellcolor{lightgreen}{93} 
%            & \cellcolor{lightgreen}{96} & \cellcolor{lightgreen}{5} & \cellcolor{lightgreen}{93} \\
%     \specialrule{.1em}{.1em}{.1em}
%   \end{tabular}
% \caption{\textbf{CR (\%) / CD (cm) / AUC (\%) comparison on the OmniObject3D test set.} Both Hestia and Hestia-H3K outperform other baselines and achieve robust performance across different object positions.}
% \label{atab:base_bench_omni}
% \end{table*}

% \begin{table*}[thb!]
%   \centering
%   \setlength{\tabcolsep}{3.1pt}
% \begin{tabular}{l|ccc|ccc|ccc|ccc|ccc}
%      & \multicolumn{15}{c}{Objaverse Test} \\
%     \specialrule{.1em}{.1em}{.1em}
%      \multirow{2}{*}{Method} 
%      & \multicolumn{3}{c|}{(0, 0)} 
%      & \multicolumn{3}{c|}{(4, 4)} 
%      & \multicolumn{3}{c|}{(4, -4)} 
%      & \multicolumn{3}{c|}{(-4, 4)} 
%      & \multicolumn{3}{c}{(-4, -4)} \\
%      & CR$\uparrow$ & CD$\downarrow$ & AUC$\uparrow$
%      & CR$\uparrow$ & CD$\downarrow$ & AUC$\uparrow$
%      & CR$\uparrow$ & CD$\downarrow$ & AUC$\uparrow$
%      & CR$\uparrow$ & CD$\downarrow$ & AUC$\uparrow$
%      & CR$\uparrow$ & CD$\downarrow$ & AUC$\uparrow$ \\
%     \specialrule{.1em}{.1em}{.1em}
%      NeU-NBV~\cite{jin2023neu} & 88 & 20 & 83 & 76 & 41 & 74 & 77 & 31 & 72 & 77 & 41 & 74 & 78 & 31 & 72 \\
%     \specialrule{.1em}{.1em}{.1em}
%      ScanRL~\cite{peralta2020next} & 87 & 18 & 82 & 80 & 31 & 78 & 82 & 23 & 77 & 77 & 37 & 76 & 80 & 27 & 76 \\
%     MACARONS~\cite{guedon2023macarons} & 88 & 17 & 79 & 91 & 16 & 86 & \cellcolor{lightyellow}{91} & \cellcolor{lightyellow}{15} & \cellcolor{lightyellow}{83} & 75 & 42 & 70 & 78 & 30 & 64 \\
%     GenNBV~\cite{chen2024gennbv} & \cellcolor{lightyellow}{94} & \cellcolor{lightyellow}{11} & \cellcolor{lightyellow}{89} & 90 & 15 & 85 & 89 & 17 & 82 & \cellcolor{lightyellow}{93} & \cellcolor{lightyellow}{12} & \cellcolor{lightyellow}{89} & \cellcolor{lightyellow}{91} & \cellcolor{lightyellow}{13} & \cellcolor{lightyellow}{86} \\
%     GenNBV (Rep.)~\cite{chen2024gennbv} & \cellcolor{lightyellow}{94} & \cellcolor{lightyellow}{11} & 88 & \cellcolor{lightyellow}{92} & \cellcolor{lightyellow}{12} & \cellcolor{lightyellow}{88} & 90 & \cellcolor{lightyellow}{15} & 82 & 92 & \cellcolor{lightyellow}{12} & 88 & 90 & 14 & 84 \\
%     \specialrule{.1em}{.1em}{.1em}
%     Hestia (Ours) & \cellcolor{lightgreen}{96} & \cellcolor{lightgreen}{7} & \cellcolor{lightgreen}{93} 
%            & \cellcolor{lightgreen}{96} & \cellcolor{lightgreen}{7} & \cellcolor{lightgreen}{92} 
%            & \cellcolor{lightgreen}{96} & \cellcolor{lightgreen}{6} & \cellcolor{lightgreen}{91} 
%            & \cellcolor{lightgreen}{96} & \cellcolor{lightgreen}{7} & \cellcolor{lightgreen}{93} 
%            & \cellcolor{lightgreen}{96} & \cellcolor{lightgreen}{8} & \cellcolor{lightgreen}{91} \\
%     \specialrule{.1em}{.1em}{.1em}
%   \end{tabular}
% \caption{\textbf{CR (\%) / CD (cm) / AUC (\%) comparison on the Objaverse test set.} Both Hestia and Hestia-H3K outperform other baselines and achieve robust performance across different object positions.}
% \label{atab:base_bench_objaverse}
% \end{table*}

% \begin{table*}[thb!]
%   \centering
%   \setlength{\tabcolsep}{3.1pt}
% \begin{tabular}{l|
%   ccc|ccc|ccc|ccc|ccc}
%      & \multicolumn{15}{c}{Houses3K Test} \\
%     \specialrule{.1em}{.1em}{.1em}
%      \multirow{2}{*}{Method} 
%      & \multicolumn{3}{c|}{(0, 0)} 
%      & \multicolumn{3}{c|}{(4, 4)} 
%      & \multicolumn{3}{c|}{(4, -4)} 
%      & \multicolumn{3}{c|}{(-4, 4)} 
%      & \multicolumn{3}{c}{(-4, -4)} \\
%      & CR$\uparrow$ & CD$\downarrow$ & AUC$\uparrow$
%      & CR$\uparrow$ & CD$\downarrow$ & AUC$\uparrow$
%      & CR$\uparrow$ & CD$\downarrow$ & AUC$\uparrow$
%      & CR$\uparrow$ & CD$\downarrow$ & AUC$\uparrow$
%      & CR$\uparrow$ & CD$\downarrow$ & AUC$\uparrow$ \\
%     \specialrule{.1em}{.1em}{.1em}
%      NeU-NBV~\cite{jin2023neu} & 85 & 26 & 80 & 73 & 52 & 71 & 81 & 30 & 77 & 72 & 53 & 70 & 80 & 32 & 76 \\
%     \specialrule{.1em}{.1em}{.1em}
%      ScanRL~\cite{peralta2020next} & 86 & 21 & 79 & 77 & 40 & 75 & 88 & 19 & 82 & 72 & 50 & 69 & 81 & 28 & 76 \\
%     MACARONS~\cite{guedon2023macarons} & 87 & 20 & 78 & 92 & 16 & 87 & 93 & 14 & 85 & 71 & 51 & 66 & 81 & 30 & 68 \\
%     GenNBV~\cite{chen2024gennbv} & \cellcolor{lightyellow}{94} & \cellcolor{lightyellow}{10} & \cellcolor{lightyellow}{89} & 90 & 18 & 83 & 91 & 16 & 84 & \cellcolor{lightyellow}{94} & 12 & \cellcolor{lightyellow}{89} & \cellcolor{lightyellow}{93} & \cellcolor{lightyellow}{12} & \cellcolor{lightyellow}{89} \\
%     GenNBV (Rep.)~\cite{chen2024gennbv} & \cellcolor{lightyellow}{94} & 12 & 88 
%                                          & \cellcolor{lightyellow}{95} & \cellcolor{lightyellow}{10} & \cellcolor{lightyellow}{90} 
%                                          & \cellcolor{lightyellow}{94} & \cellcolor{lightyellow}{13} & \cellcolor{lightyellow}{88} 
%                                          & \cellcolor{lightyellow}{94} & \cellcolor{lightyellow}{11} & \cellcolor{lightyellow}{89} 
%                                          & 91 & 14 & 88 \\
%     \specialrule{.1em}{.1em}{.1em}
%     Hestia (Ours) & \cellcolor{lightgreen}{97} & \cellcolor{lightgreen}{5} & \cellcolor{lightgreen}{96} 
%            & \cellcolor{lightgreen}{97} & \cellcolor{lightgreen}{5} & \cellcolor{lightgreen}{93} 
%            & \cellcolor{lightgreen}{97} & \cellcolor{lightgreen}{7} & \cellcolor{lightgreen}{93} 
%            & \cellcolor{lightgreen}{97} & \cellcolor{lightgreen}{5} & \cellcolor{lightgreen}{94} 
%            & \cellcolor{lightgreen}{97} & \cellcolor{lightgreen}{7} & \cellcolor{lightgreen}{93} \\
%     \specialrule{.1em}{.1em}{.1em}
%   \end{tabular}
% \caption{\textbf{CR (\%) / CD (cm) / AUC (\%) comparison on the Houses3K test set.} Hestia consistently achieves the best performance across different object positions, while Hestia-H3K (not shown here) also outperforms prior baselines.}
% \label{atab:base_bench_h3k}
% \end{table*}

\subsection{Limited Acquisitions}

This section demonstrates the suitability of Hestia for efficient 3D reconstruction, where only a limited number of views can be acquired. 
\cref{tab:limited_acq} shows that Hestia outperforms prior works by at least 12\% and 5\% in CR under 5-image and 15-image budgets, respectively. 
Specifically, Hestia achieves 92\% CR with only 5 acquisitions, whereas prior work~\cite{chen2024gennbv} requires 15 images to reach comparable performance. 
These gains stem from the close-greedy training strategy, which not only mitigates spurious correlations over time but also enables efficient capture during inference.

% \begin{figure*}[tbh!]
%   \centering
%   \includegraphics[width=\linewidth]{figures/curves.png}
%   \caption{\textbf{CR curves on three datasets.} Hestia outperforms prior approaches by nearly 10\% and 5\% in the first five captures and the first fifteen captures, respectively. The efficiency is particularly significant in real-world power-constrained scenarios, where a robot or agent may run out of battery in a short time.}
%   \label{fig:fail}
% \end{figure*}

\begin{table}[thb!]
\centering
\setlength{\tabcolsep}{3.1pt}
\begin{tabular}{l|cc}
Method & 5 images & 15 images \\
    \specialrule{.1em}{.1em}{.1em}
    NeU-NBV~\cite{jin2023neu} & 71 & 76\\
    ScanRL~\cite{peralta2020next} & 69 & 80\\
    MACARONS~\cite{guedon2023macarons} & 63 & 82\\
    GenNBV~\cite{chen2024gennbv} & 75 & \cellcolor{lightyellow}{91}\\
    GenNBV (Rep.)~\cite{chen2024gennbv} & \cellcolor{lightyellow}{80} & \cellcolor{lightyellow}{91}\\
    \specialrule{.1em}{.1em}{.1em}
    Hestia (Ours) & \cellcolor{lightgreen}{92} & \cellcolor{lightgreen}{96}\\
  \end{tabular}
\caption{\textbf{Mean CR (\%) comparison across the three datasets with limited-view acquisition.} Hestia outperforms prior approaches by at least 12\% and 5\% with a 5-image budget and a 15-image budget. Such efficiency is crucial in real-world power-constrained settings, as robots or agents may exhaust their battery within a short time.}
\label{tab:limited_acq}
\end{table}

\subsection{Inference Speed}

Inference speed is critical for next-best-view planning because robots with onboard cameras must capture images for 3D reconstruction before their batteries are depleted. 
As shown in \cref{tab:avg_bench}, Hestia achieves 25 FPS, which is suitable for real-time deployment. 
Modeling voxels as cubes rather than points does not significantly reduce inference speed. 
The speed also demonstrates the advantage of RL-based generalizable next-best-view approaches, since using a policy model to directly predict viewpoints removes the need to sample candidate views for prediction.

\subsection{Ablation Study}
\label{sec:abl}

This section evaluates the effectiveness of Hestia’s core components through an ablation study (see~\cref{tab:ablation}). 
We investigate three key ideas: face-aware design, a close-greedy training scheme, and a hierarchical structure. 
For the non-hierarchical variant, the encoded grid information is fed directly into the policy model to predict 5-DoF viewpoints without feature interpolation, since interpolation requires the look-at point. 
Applying the close-greedy strategy or the hierarchical structure alone eases the training process, whereas face-aware observation alone may make training more difficult but still provides complementary information. 
Thus, the close-greedy strategy and hierarchical structure yield stronger gains when applied individually, while combining them with face-aware observation further enhances stability and capture quality. 
Overall, each component contributes to performance improvements, and integrating all three delivers the best results.

\begin{table}[thb!]
  \centering
\begin{tabular}{@{}l|ccc|cc|c}
    & Face & Greedy & Hier. & CR $\uparrow$ & CD $\downarrow$ & \#Pa. \\
    \specialrule{.1em}{.1em}{.1em}
    \multirow{9}{*}{Hestia} &  & & & 88 & 20 & 6.2M \\
    \cmidrule{2-7}
    & \checkmark & & & 90 & 17 & 6.2M \\
    & & \checkmark & & 92 & 13  & 6.2M \\
    & & & \checkmark & 94 & 11 & 4.9M \\
    \cmidrule{2-7}
    & \checkmark & \checkmark & & 95 & 9 & 6.2M \\
    & & \checkmark & \checkmark & 95 & 8 & 4.9M \\
    & \checkmark & & \checkmark & 95 & 9 & 4.9M \\
    \cmidrule{2-7}
    & \checkmark & \checkmark & \checkmark & 96 & 7 & 4.9M\\
    \specialrule{.1em}{.1em}{.1em}
  \end{tabular}
\caption{\textbf{Ablations on Objaverse~\cite{deitke2024objaverse, Deitke_2023_CVPR}.} Integrating the proposed ideas yields the best performance with fewer parameters.}
  \label{tab:ablation}
\end{table}

\subsection{Application}

\begin{figure*}[tbh!]
  \centering
  \includegraphics[width=0.82\linewidth]{figures/real_world_vis.png}
  \caption{\textbf{Real-world demonstration of non-shifted and shifted scenes.} Red boxes indicate manually initialized viewpoints, while blue boxes denote viewpoints predicted by Hestia. The results demonstrate Hestia’s feasibility in real-world environments.}
  \label{fig:vis_real}
\end{figure*}

% \begin{figure*}[tbh!]
%   \centering
%   \includegraphics[width=\linewidth]{figures/real_world_vis.png}
%   \caption{\textbf{Real-world demonstration of non-shifted and shifted scenes.} Red boxes indicate manually initialized viewpoints, while blue boxes denote viewpoints predicted by Hestia. The results demonstrate Hestia’s feasibility in real-world environments.}
%   \label{fig:vis_real}
% \end{figure*}

This section demonstrates that Hestia is feasible for real-world deployment even without a depth camera. 
We use a drone equipped with an RGB camera as the mobile agent for data collection and employ a depth predictor~\cite{duisterhof2024mast3r, wang2024dust3r} to convert multi-view RGB images into depth maps. 
The first three images are manually selected to synchronize the real-world and virtual-world settings. 
As shown in \cref{fig:vis_real}, Hestia successfully operates in real-world scenarios for both shifted and non-shifted scenes.
It is worth noting that some prior works~\cite{peralta2020next, chen2024gennbv, jin2023neu, guedon2023macarons} report only simulation results, and their real-world feasibility remains unknown. 
Please see~\cref{asec:realworld_res,asec:realworld_sys} for more details.



